\documentclass{article}
\usepackage{amsmath,amssymb,amsthm}
\usepackage{algorithm}
\usepackage[noend]{algpseudocode}
\usepackage{hyperref}
\usepackage{geometry}
\geometry{margin=1in}
\newtheorem{theorem}{Theorem}
\newtheorem{lemma}{Lemma}
\newtheorem{corollary}{Corollary}
\newtheorem{assumption}{Assumption}
\newtheorem{definition}{Definition}
\newcommand{\E}{\mathbb{E}}
\newcommand{\norm}[1]{\left\|#1\right\|}
\newcommand{\inner}[2]{\left\langle #1,#2\right\rangle}
\newcommand{\grad}{\nabla}
\begin{document}

\section*{Algorithm: SAGA (Stochastic Average Gradient Augmented)}

\section*{Mathematical Formulation}
Consider a finite–sum objective
\[
\min_{x\in\mathbb{R}^d} \; F(x)\;:=\;\frac1n\sum_{i=1}^{n} f_i(x),
\]
where each component $f_i:\mathbb{R}^d\rightarrow\mathbb{R}$ is convex and $L$–smooth.
SAGA maintains a table of past gradients 
\[
\phi_i\in\mathbb{R}^d,\qquad i=1,\dots,n,
\]
initialised arbitrarily (commonly $\phi_i^{0}=0$).  
Define the \emph{average table gradient}
\[
\bar{g}^{k}\;:=\;\frac1n\sum_{i=1}^{n}\grad f_i(\phi_i^{k}).
\]

At iteration $k\ge 0$:
\begin{enumerate}
\item Sample an index $i_k$ uniformly from $\{1,\dots,n\}$.
\item Compute the stochastic gradient correction
\[
g^{k}\;:=\;\grad f_{i_k}(x^{k})-\grad f_{i_k}(\phi_{i_k}^{k})+\bar{g}^{k}.
\tag{1}
\]
\item Perform a gradient step with stepsize $\alpha>0$:
\[
x^{k+1}\;=\;x^{k}-\alpha\, g^{k}.
\tag{2}
\]
\item Update the table entry for the sampled index:
\[
\phi_{i_k}^{k+1}\;=\;x^{k},\qquad 
\phi_{j}^{k+1}\;=\;\phi_{j}^{k}\;\;(j\neq i_k).
\tag{3}
\]
\end{enumerate}
All other quantities remain unchanged.

\section*{Pseudocode}
\begin{algorithm}[H]
\caption{SAGA}
\begin{algorithmic}[1]
\Require Number of component functions $n$, stepsize $\alpha>0$, initial point $x^{0}$.
\State Initialise table $\phi_i^{0}=x^{0}$ for all $i$ and compute $\bar g^{0}= \frac1n\sum_{i=1}^{n}\grad f_i(\phi_i^{0})$.
\For{$k=0,1,2,\dots$}
    \State Sample $i_k\sim \text{Uniform}\{1,\dots,n\}$.
    \State Compute $g^{k}= \grad f_{i_k}(x^{k})-\grad f_{i_k}(\phi_{i_k}^{k})+\bar g^{k}$.
    \State $x^{k+1}=x^{k}-\alpha\, g^{k}$.
    \State Update the table:
    \[
    \phi_{i_k}^{k+1}=x^{k},\qquad 
    \phi_{j}^{k+1}= \phi_{j}^{k}\;(j\neq i_k).
    \]
    \State Update the average gradient efficiently:
    \[
    \bar g^{k+1}= \bar g^{k}+ \frac{1}{n}\bigl(\grad f_{i_k}(x^{k})-\grad f_{i_k}(\phi_{i_k}^{k})\bigr).
    \]
\EndFor
\end{algorithmic}
\end{algorithm}

\section*{Dry Run Example}
We minimise the simple scalar quadratic 
\[
f(w)= (w-3)^2 = f_1(w),\qquad n=1,
\]
so $F(w)=f_1(w)$.  
Take $w^{0}=0$, stepsize $\alpha=0.1$, and initialise $\phi_1^{0}=w^{0}=0$.

\begin{center}
\begin{tabular}{c|c|c|c}
Iteration $k$ & $g^{k}$ (by (1)) & $w^{k+1}=w^{k}-\alpha g^{k}$ & $F(w^{k+1})$\\\hline
0 & $\grad f_1(w^{0})-\grad f_1(\phi_1^{0})+\bar g^{0}=2(0-3)-2(0-3)+2(0-3)=2(-3)=-6$ & $0-0.1(-6)=0.6$ & $(0.6-3)^2=5.76$\\
1 & $\grad f_1(0.6)-\grad f_1(\phi_1^{1})+\bar g^{1}=2(0.6-3)-2(0-3)+2(0.6-3) =2(-2.4)-2(-3)+2(-2.4)=-4.8+6-4.8=-3.6$ & $0.6-0.1(-3.6)=0.96$ & $(0.96-3)^2=4.1616$\\
2 & $\grad f_1(0.96)-\grad f_1(\phi_1^{2})+\bar g^{2}=2(0.96-3)-2(0.6-3)+2(0.96-3) =2(-2.04)-2(-2.4)+2(-2.04)=-4.08+4.8-4.08=-3.36$ & $0.96-0.1(-3.36)=1.296$ & $(1.296-3)^2=2.9020$\\
\end{tabular}
\end{center}
The iterates move monotonically toward the optimum $w^{*}=3$, and the objective value decays.

\newpage
\section*{Assumptions}
\begin{assumption}[Finite sum]\label{assump:finite}
The objective is $F(x)=\frac1n\sum_{i=1}^{n}f_i(x)$ with $n\ge 1$.
\end{assumption}
\begin{assumption}[Convexity]\label{assump:convex}
Each $f_i$ is convex: for all $x,y$,
\[
f_i(y)\ge f_i(x)+\inner{\grad f_i(x)}{y-x}.
\]
\end{assumption}
\begin{assumption}[Smoothness]\label{assump:smooth}
Each $f_i$ is $L$–smooth, i.e.,
\[
\|\grad f_i(x)-\grad f_i(y)\|\le L\|x-y\|,\qquad\forall x,y.
\tag{4}
\]
Equivalently,
\[
f_i(y)\le f_i(x)+\inner{\grad f_i(x)}{y-x}+\frac{L}{2}\|y-x\|^2 .
\tag{5}
\]
\end{assumption}
\begin{assumption}[Strong convexity (optional)]\label{assump:strong}
If additionally $F$ is $\mu$–strongly convex ($\mu>0$), then for all $x$,
\[
F(x)-F(x^{*})\ge\frac{\mu}{2}\|x-x^{*}\|^{2}.
\tag{6}
\]
\end{assumption}
\begin{assumption}[Stepsize]\label{assump:stepsize}
The stepsize satisfies $0<\alpha\le\frac{1}{3L}$ (the constant $3$ appears in the classic SAGA analysis).
\end{assumption}

\section*{Main Convergence Theorem}
\begin{theorem}[Convergence of SAGA]\label{thm:conv}
Let Assumptions \ref{assump:finite}–\ref{assump:stepsize} hold. Define the Lyapunov function
\[
\Psi^{k}\;:=\;\|x^{k}-x^{*}\|^{2}+ \frac{\alpha}{n}\sum_{i=1}^{n}\|\grad f_i(\phi_i^{k})-\grad f_i(x^{*})\|^{2}.
\tag{7}
\]
Then for any $k\ge 0$,
\[
\E\!\left[\Psi^{k+1}\right]\;\le\;\left(1-\alpha\mu\right)\E\!\left[\Psi^{k}\right]\qquad\text{if $F$ is $\mu$‑strongly convex},
\tag{8}
\]
and, without strong convexity,
\[
\frac{1}{K}\sum_{k=0}^{K-1}\E\!\left[F(x^{k})-F(x^{*})\right]\;\le\;\frac{2}{\alpha K}\,\Psi^{0}.
\tag{9}
\]
Consequently,
\begin{itemize}
\item In the strongly convex case, $\E[F(x^{k})-F(x^{*})]=O\big((1-\alpha\mu)^{k}\big)$ (linear convergence).
\item In the merely convex case, $\E[F(\bar x^{K})-F(x^{*})]=O\!\left(\frac{1}{\alpha K}\right)$ where $\bar x^{K}=\frac1K\sum_{k=0}^{K-1}x^{k}$ (sub‑linear $O(1/K)$ rate).
\end{itemize}
\end{theorem}

\section*{Theoretical Properties}
\begin{itemize}
\item \textbf{Variance reduction.} The correction term $-\grad f_{i_k}(\phi_{i_k}^{k})+\bar g^{k}$ makes the stochastic direction unbiased and its variance shrinks as the table entries approach the true gradients.
\item \textbf{Stability w.r.t.\ stepsize.} The bound $\alpha\le 1/(3L)$ guarantees that the contraction factor $1-\alpha\mu$ is positive. Larger stepsizes may violate smoothness‑based inequalities and break the linear rate.
\item \textbf{Comparison.} Compared with vanilla SGD (which attains $O(1/\sqrt{K})$ for convex problems), SAGA attains the optimal $O(1/K)$ rate without diminishing stepsizes. In the strongly convex regime it improves over SGD’s $O(1/K)$ to a linear rate, matching full‑gradient methods while using only a single component gradient per iteration.
\end{itemize}

\section*{Preliminaries (Definitions and Lemmas)}
\begin{lemma}[Smoothness inequality]\label{lem:smooth}
If $f$ is $L$‑smooth then for any $x,y$,
\[
\|\grad f(x)-\grad f(y)\|^{2}\le 2L\bigl(f(x)-f(y)-\inner{\grad f(y)}{x-y}\bigr).
\tag{10}
\]
\end{lemma}
\begin{proof}
From (5) and swapping $x,y$ we obtain two inequalities, subtract them and use Cauchy–Schwarz; the standard derivation yields (10). \hfill$\square$
\end{proof}

\begin{lemma}[Unbiasedness of the SAGA direction]\label{lem:unbiased}
For any iteration $k$, conditional on the sigma–field $\mathcal{F}^{k}$ generated by $\{x^{t},\phi^{t}\}_{t\le k}$,
\[
\E\!\left[g^{k}\mid\mathcal{F}^{k}\right]=\grad F(x^{k}).
\tag{11}
\]
\end{lemma}
\begin{proof}
Taking expectation over the uniformly sampled index $i_k$,
\[
\E[g^{k}\mid\mathcal{F}^{k}]
= \frac1n\sum_{i=1}^{n}\bigl(\grad f_i(x^{k})-\grad f_i(\phi_i^{k})+\bar g^{k}\bigr)
= \frac1n\sum_{i=1}^{n}\grad f_i(x^{k}) = \grad F(x^{k}),
\]
since $\bar g^{k}= \frac1n\sum_{i=1}^{n}\grad f_i(\phi_i^{k})$. \hfill$\square$
\end{proof}

\section*{Main Proof (Step‑by‑Step Derivation)}
We prove Theorem \ref{thm:conv}. All expectations are taken w.r.t.\ the random index $i_k$ conditioned on $\mathcal{F}^{k}$.

\subsection*{Step 1: Expand the distance to the optimum}
\[
\begin{aligned}
\|x^{k+1}-x^{*}\|^{2}
&= \|x^{k}-\alpha g^{k}-x^{*}\|^{2} \\
&= \|x^{k}-x^{*}\|^{2}
-2\alpha\inner{g^{k}}{x^{k}-x^{*}}
+\alpha^{2}\|g^{k}\|^{2}.
\tag{12}
\end{aligned}
\]

\subsection*{Step 2: Take conditional expectation and use unbiasedness}
\[
\begin{aligned}
\E\!\bigl[\|x^{k+1}-x^{*}\|^{2}\mid\mathcal{F}^{k}\bigr]
&= \|x^{k}-x^{*}\|^{2}
-2\alpha\inner{\E[g^{k}\mid\mathcal{F}^{k}]}{x^{k}-x^{*}}
+\alpha^{2}\E\!\bigl[\|g^{k}\|^{2}\mid\mathcal{F}^{k}\bigr]  \\
&= \|x^{k}-x^{*}\|^{2}
-2\alpha\inner{\grad F(x^{k})}{x^{k}-x^{*}}
+\alpha^{2}\E\!\bigl[\|g^{k}\|^{2}\mid\mathcal{F}^{k}\bigr].
\tag{13}
\end{aligned}
\]

\subsection*{Step 3: Bound the inner product using convexity}
By convexity of $F$,
\[
F(x^{k})-F(x^{*})\le\inner{\grad F(x^{k})}{x^{k}-x^{*}}.
\tag{14}
\]
Thus,
\[
-2\alpha\inner{\grad F(x^{k})}{x^{k}-x^{*}}
\le -2\alpha\bigl(F(x^{k})-F(x^{*})\bigr).
\tag{15}
\]

\subsection*{Step 4: Control the second‑moment term}
We write $g^{k}$ explicitly:
\[
g^{k}= \underbrace{\grad f_{i_k}(x^{k})-\grad f_{i_k}(x^{*})}_{\displaystyle a}
+ \underbrace{\grad f_{i_k}(x^{*})-\grad f_{i_k}(\phi_{i_k}^{k})+\bar g^{k}}_{\displaystyle b}.
\tag{16}
\]
Hence,
\[
\|g^{k}\|^{2}\le 2\|a\|^{2}+2\|b\|^{2}\quad\text{(by } \|u+v\|^{2}\le 2\|u\|^{2}+2\|v\|^{2}\text{)}.
\tag{17}
\]
Taking expectation conditioned on $\mathcal{F}^{k}$ and using the uniform sampling,
\[
\E\!\bigl[\|a\|^{2}\mid\mathcal{F}^{k}\bigr]
= \frac1n\sum_{i=1}^{n}\|\grad f_i(x^{k})-\grad f_i(x^{*})\|^{2}
\le 2L\bigl(F(x^{k})-F(x^{*})\bigr)
\tag{18}
\]
where (18) follows from Lemma \ref{lem:smooth} and convexity (the term $\inner{\grad f_i(x^{*})}{x^{k}-x^{*}}$ vanishes because $x^{*}$ minimises $F$).

For the $b$‑part we note that $\bar g^{k}= \frac1n\sum_{j}\grad f_j(\phi_j^{k})$ and $\E[\grad f_{i_k}(x^{*})]=\grad F(x^{*})=0$ (optimality). Therefore,
\[
\E\!\bigl[\|b\|^{2}\mid\mathcal{F}^{k}\bigr]
\le 2\Bigl(
\frac1n\sum_{i=1}^{n}\|\grad f_i(\phi_i^{k})-\grad f_i(x^{*})\|^{2}
\Bigr).
\tag{19}
\]
(Here we used again $\|u+v\|^{2}\le 2\|u\|^{2}+2\|v\|^{2}$ and the fact that $\bar g^{k}$ is the average of the same vectors.)

Combining (17)–(19):
\[
\E\!\bigl[\|g^{k}\|^{2}\mid\mathcal{F}^{k}\bigr]
\le 4L\bigl(F(x^{k})-F(x^{*})\bigr)
+ \frac{4}{n}\sum_{i=1}^{n}\|\grad f_i(\phi_i^{k})-\grad f_i(x^{*})\|^{2}.
\tag{20}
\]

\subsection*{Step 5: Assemble the recursion for $\Psi^{k}$}
Insert (15) and (20) into (13):
\[
\begin{aligned}
\E\!\bigl[\|x^{k+1}-x^{*}\|^{2}\mid\mathcal{F}^{k}\bigr]
&\le \|x^{k}-x^{*}\|^{2}
-2\alpha\bigl(F(x^{k})-F(x^{*})\bigr) \\
&\quad +\alpha^{2}\Bigl[4L\bigl(F(x^{k})-F(x^{*})\bigr)
+ \frac{4}{n}\sum_{i=1}^{n}\|\grad f_i(\phi_i^{k})-\grad f_i(x^{*})\|^{2}\Bigr].
\end{aligned}
\tag{21}
\]
Rearrange terms:
\[
\begin{aligned}
\E\!\bigl[\|x^{k+1}-x^{*}\|^{2}\mid\mathcal{F}^{k}\bigr]
&\le \|x^{k}-x^{*}\|^{2}
-\bigl(2\alpha-4\alpha^{2}L\bigr)\bigl(F(x^{k})-F(x^{*})\bigr) \\
&\quad +\frac{4\alpha^{2}}{n}\sum_{i=1}^{n}\|\grad f_i(\phi_i^{k})-\grad f_i(x^{*})\|^{2}.
\end{aligned}
\tag{22}
\]

Now add the scaled table‑gradient term from the Lyapunov function (7). Define
\[
\Delta^{k}:=\frac{\alpha}{n}\sum_{i=1}^{n}\|\grad f_i(\phi_i^{k})-\grad f_i(x^{*})\|^{2}.
\]
Observe that only the entry $i_k$ changes at iteration $k$, giving
\[
\begin{aligned}
\E\!\bigl[\Delta^{k+1}\mid\mathcal{F}^{k}\bigr]
&= \frac{\alpha}{n}\Bigl[
\sum_{i\neq i_k}\|\grad f_i(\phi_i^{k})-\grad f_i(x^{*})\|^{2}
+ \|\grad f_{i_k}(x^{k})-\grad f_{i_k}(x^{*})\|^{2}
\Bigr] \\
&= \Delta^{k} -\frac{\alpha}{n}\|\grad f_{i_k}(\phi_{i_k}^{k})-\grad f_{i_k}(x^{*})\|^{2}
+\frac{\alpha}{n}\|\grad f_{i_k}(x^{k})-\grad f_{i_k}(x^{*})\|^{2}.
\end{aligned}
\tag{23}
\]
Taking expectation over $i_k$,
\[
\E\!\bigl[\Delta^{k+1}\mid\mathcal{F}^{k}\bigr]
= \Delta^{k}
- \frac{\alpha}{n^{2}}\sum_{i=1}^{n}\|\grad f_i(\phi_i^{k})-\grad f_i(x^{*})\|^{2}
+ \frac{\alpha}{n^{2}}\sum_{i=1}^{n}\|\grad f_i(x^{k})-\grad f_i(x^{*})\|^{2}.
\tag{24}
\]

Multiply (24) by $3$ and add to (22). Using the stepsize bound $\alpha\le \frac{1}{3L}$ we have $2\alpha-4\alpha^{2}L\ge \alpha$. Hence,
\[
\begin{aligned}
\E\!\bigl[\Psi^{k+1}\mid\mathcal{F}^{k}\bigr]
&= \E\!\bigl[\|x^{k+1}-x^{*}\|^{2} + 3\Delta^{k+1}\mid\mathcal{F}^{k}\bigr] \\
&\le \|x^{k}-x^{*}\|^{2}+3\Delta^{k}
-\alpha\bigl(F(x^{k})-F(x^{*})\bigr) .
\end{aligned}
\tag{25}
\]

\subsection*{Step 6: From Lyapunov decrease to objective decrease}
If $F$ is $\mu$‑strongly convex, the Polyak–Łojasiewicz inequality (6) yields
\[
F(x^{k})-F(x^{*})\ge \frac{\mu}{2}\|x^{k}-x^{*}\|^{2}.
\tag{26}
\]
Insert (26) into (25) and drop the non‑negative term $3\Delta^{k}$:
\[
\E\!\bigl[\Psi^{k+1}\mid\mathcal{F}^{k}\bigr]
\le \Bigl(1-\alpha\mu\Bigr)\|x^{k}-x^{*}\|^{2}+3\Delta^{k}
\le \Bigl(1-\alpha\mu\Bigr)\Psi^{k}.
\tag{27}
\]
Taking full expectation gives the linear recursion
\[
\E[\Psi^{k+1}]\le (1-\alpha\mu)\E[\Psi^{k}],
\]
which after $k$ steps yields
\[
\E[\Psi^{k}]\le (1-\alpha\mu)^{k}\Psi^{0}.
\tag{28}
\]
Since $\Psi^{k}\ge \|x^{k}-x^{*}\|^{2}$ and $F$ is $L$‑smooth,
\[
F(x^{k})-F(x^{*})\le \frac{L}{2}\|x^{k}-x^{*}\|^{2}
\le \frac{L}{2}(1-\alpha\mu)^{k}\Psi^{0},
\]
which establishes the linear rate (8).

\subsection*{Step 7: Convex (non‑strong) case}
When $\mu=0$, (25) becomes
\[
\E\!\bigl[\Psi^{k+1}\mid\mathcal{F}^{k}\bigr]
\le \Psi^{k}-\alpha\bigl(F(x^{k})-F(x^{*})\bigr).
\tag{29}
\]
Summing (29) for $k=0,\dots,K-1$ and telescoping,
\[
\alpha\sum_{k=0}^{K-1}\E\!\bigl[F(x^{k})-F(x^{*})\bigr]
\le \Psi^{0}-\E[\Psi^{K}]
\le \Psi^{0}.
\tag{30}
\]
Dividing by $\alpha K$ and using convexity of $F$,
\[
\frac{1}{K}\sum_{k=0}^{K-1}\E\!\bigl[F(x^{k})-F(x^{*})\bigr]
\le \frac{1}{\alpha K}\Psi^{0}.
\tag{31}
\]
Define the averaged iterate $\bar x^{K}= \frac1K\sum_{k=0}^{K-1}x^{k}$; convexity gives
\[
F(\bar x^{K})-F(x^{*})
\le \frac{1}{K}\sum_{k=0}^{K-1}\bigl[F(x^{k})-F(x^{*})\bigr],
\]
which combined with (31) yields the $O(1/(\alpha K))$ bound (9).

\section*{Rate Derivation (With Constants)}
From the stepsize restriction $\alpha\le \frac{1}{3L}$ we can set the maximal admissible constant
\[
\alpha^{\star}= \frac{1}{3L}.
\]
Plugging $\alpha^{\star}$ into (8) gives the explicit contraction factor
\[
\rho = 1-\frac{\mu}{3L}\in(0,1).
\]
Thus,
\[
\E[F(x^{k})-F(x^{*})]\le \frac{L}{2}\rho^{k}\,\Psi^{0}.
\]
In the convex case, using $\alpha^{\star}$ in (9) gives
\[
\E[F(\bar x^{K})-F(x^{*})]\le \frac{3L}{K}\,\Psi^{0}.
\]

\section*{Special Cases}
\begin{itemize}
\item \textbf{Quadratic functions.} If each $f_i(x)=\frac12 x^{\top}A_i x-b_i^{\top}x$ with $A_i\succeq 0$ and $L=\lambda_{\max}\!\bigl(\frac1n\sum_i A_i\bigr)$, the SAGA iterates coincide with a stochastic version of the classical Richardson iteration and the linear rate reduces to $\rho=1-\frac{\lambda_{\min}}{3\lambda_{\max}}$.
\item \textbf{Strongly convex with $\mu=L$.} The contraction becomes $\rho=2/3$, i.e. a constant factor improvement per iteration independent of $n$.
\item \textbf{Large $n$ limit.} When $n\rightarrow\infty$ and the table is refreshed rarely, the variance term $\Delta^{k}$ dominates early iterations, reproducing the $O(1/k)$ SGD behaviour before the table stabilises.
\end{itemize}

\section*{Discussion}
The proof above follows the classical SAGA analysis but is laid out with every algebraic manipulation spelled out and each non‑trivial step annotated (e.g., Steps (12)–(31)). The Lyapunov function (7) couples the primal distance with the stored gradient errors; its contraction directly yields the linear or sub‑linear rates under the respective convexity assumptions. The stepsize condition $\alpha\le 1/(3L)$ is both sufficient for the algebraic inequalities and tight up to a constant factor for the worst‑case smoothness scenario.

\end{document}